\tocsection{\textbf{9/29/26 Lecture 9:} Lebesgue-Fubini-Tonelli-H\"older-Minkowski}

\begin{lemma}[\textbf{Lemma 9.1:} Completion]
    Consider a measure space $(\Omega, \cal{A}, \mu)$, a $\sigma$-field $\cal{F} \subseteq \cal{A}$, and a Borel space $(S, \cal{S})$. Then, a function $f : \Omega \to S$ is $\cal{F}^{\mu}$-measurable if, and only if, there exists some $\cal{F}$-measurable function $g$ such that $f = g$ a.e. $\mu$. 
\end{lemma}
\begin{proof}{Proof:}
    Let's start with indicators. If $\cal{G}$ is the set of $A \subseteq \Omega$ such that $A \Delta B \in \cal{N}_{\mu}$ for some $B \in \cal{F}$. Then, we have that $A \backslash B$ and $B \backslash A$ are in $\cal{N}_{\mu}$ (they are subset of a null set), which implies that $\cal{G} \subseteq \cal{F}^{\mu}$. This is because we know that $B \cup A \backslash B = A \cup B \in \cal{F}^{\mu}$ and $A \cup B \backslash (B \backslash A) = A \in \cal{F}^{\mu}$. Conversely, $\cal{F}^{\mu} \subseteq \cal{G}$ since both $\cal{F}$ and $\cal{N}_{\mu}$ are trivially ($\cal{F}$ is trivial, for $\cal{N}_{\mu}$, consider $A \Delta \emptyset$) contained in $\cal{G}$, and it can be shown that $\cal{G}$ is a $\sigma$-field. \\

    Then, we see that $A \in \cal{F}^{\mu}$ if, and only if $1_{A} = 1_{B}$ almost everywhere for some $B \in \cal{F}$. When, we may choose some simple $\cal{F}^{\mu}$-measurable functions $f_n$ such that $0 \leq f_n \uparrow f$. By our result above there exists $\cal{F}$-measurable functions $g_n$ such that $f_n = g_n$ almost everywhere. Since a countable union of null sets is a null set, we see that $g = \limsup_{n} g_n$ has the property that we want.

    $\qed$
\end{proof}

\begin{definition}[\textbf{Definition 9.2:} $\sigma$-finiteness]
    A measure $\mu$ on a measurable space $(\Omega, \cal{A})$ is \emph{$\sigma$-finite} if $\Omega = \bigcup_{i = 1}^{\infty} A_n$ with $A_i \cap A_j = \emptyset$ for $i \ne j$ such that $\mu A_n < \infty$ for all $n$. Equivalently, $\mu$ is $\sigma$-finite if, and only if, there is a measurable strictly positive function $f$ such that $\mu f < \infty$. 
\end{definition}


The following results answer when
\[ \int_{S}\left(\int_{T} f(s, t) \nu(t)\right)\mu(s) = \int_{T}\left(\int_{S} f(s, t)\mu(s)\right)\mu(t) \]
First, we start with a lemma to show that the statement above even makes sense at all.
\begin{lemma}[\textbf{Lemma 9.3:} Sections]
    Fix two measurable spaces $(S, \cal{S})$ and $(T, \cal{T})$, a measurable function $f : S \times T \to \R_{+}$, and a $\sigma$-finite measure $\mu$ on $S$. Then,
    \begin{enumerate}
        \item $f(s, t)$ is $\cal{S}$ measurable in $s \in S$ for each $t \in T$.
        \item $\int f(s, t)\mu(ds)$ is $\cal{T}$-measurable in $t \in T$. 
    \end{enumerate}
\end{lemma}
\begin{proof}{Proof:}
    We may assume that $\mu$ is bounded. Both statements are obvious when $f = 1_{A}$ with $A = B \times C$ for some $B \in \cal{S}$ and $C \in \cal{T}$. This extends by a monotone class argument to any indicator functions of sets in $\cal{S} \otimes \cal{T}$. The general case then follows by linearity and monotone convergence. 

    $\qed$
\end{proof}

\begin{theorem}[\textbf{Theorem 9.4:} Lebesgue-Fubini-Tonelli Theorem]
    For any $\sigma$-finite measure spaces $(S, \cal{S}, \mu)$ and $(T, \cal{T}, \nu)$, there exists a unique measure $\mu \otimes \nu$ on $(S \times T, \cal{S} \otimes \cal{T})$ satisfying
    \[ (\mu \otimes \nu)(B \times C) = \mu(B) \cdot \nu(C),\quad B \in \cal{S}, C \in \cal{T} \]
    Furthermore, for any non-negative measureable function $f : S \times T \to \R_{+}$, 
    \[ (\mu \otimes \nu)v = \int \mu(ds)\int f(s, t)\nu(dt) = \int \nu(dt) \int f(s, t) \mu(ds) \]
    The last relation remains valid for any measurable function $f : S \times T \to \R$ with $(\mu \otimes \nu)\abs{f} < \infty$.\\
    
    Note that these integrals are well defined by the lemma above, and although the inner integrals $\nu f(s, \cdot)$ and $\mu f(\cdot, t)$ may fail to exist on some null sets in $S$ and $T$, when everything works, everything agrees. 
\end{theorem}
\begin{proof}{Proof:}
    By the lemma above, we can define
    \[ (\mu \otimes \nu)(A) = \int \mu(ds)\int 1_{A}(s, t)\nu(dt), \quad A \in \cal{S} \otimes \cal{T} \]
    which is a clearly a measure on $S \times T$. By a monotone class argument, there can be at most one measure (Caratheodory Extension Theorem).\\

    In particular, if we interchange the integrals, from the same monotone class argument, we see that the corresponding measure $\nu \otimes \mu$ also agrees with $\mu \otimes \nu$. Thus, we see that the statement in the theorem about integrals holds for indicator functions. Then, by linearity and monotone convergence, it holds for arbitrary measurable functions $f \geq 0$.\\

    Now, for the general case, we see that for any measurable $f$, statement about integrals holds for $\abs{f}$. Then, if $(\mu \otimes \nu)\abs{f} < \infty$, it follows that $N_{S} = \{s \in S \col \nu\abs{f(s, \cdot)} = \infty\}$ is a $\mu$-null set in $S$ (as otherwise we would get infinity for our integral), and similarly $N_{T} = \{t \in T \col \mu\abs{f(\cdot, t)} = \infty\}$ is also a $\nu$-null set. Then, we can define $f(s, t)$ to  be zero when $s \in N_{S}$ or $t \in N_{t}$. This new function has the same integral as the previous function.\\

    Then, the statement about integrals follows for $f$ by subtraction of the formulas for $f_{+}$ and $f_{-}$. 

\end{proof}

Now, we move on to $L^{p}$ spaces!

\begin{definition}[\textbf{Definition 9.5:} $L^{p}$ Spaces]
    Given a measure space $(\Omega, \cal{A}, \mu)$ and a $p > 0$, write $L^{p} = L^{p}(\Omega, \cal{A}, \mu)$ for the class of all measurable functions $f : \Omega \to \R$ with
    \[ \norm{f}_{p} := \left[\int \abs{f}^{p} d\mu\right]^{1/p} < \infty \]
\end{definition}

Now, here are some very useful inequalities:

\begin{theorem}[\textbf{Theorem 9.6:} H\"older and Minkowski Inequalities]
    For any measurable functions $f$ and $g$ on some measure space $(\Omega, \cal{A}, \mu)$, we have
    \begin{enumerate}
        \item $\norm{fg}_{r} \leq \norm{f}_{p}\norm{g}_{q}$ for all $p, q, r > 0$ with $p^{-1} + q^{-1} = r^{-1}$.
        \item $\norm{f + g}_{p}\leq \norm{f}_{p} + \norm{g}_{p}$ for all $p \geq 1$, and $\norm{f + g}_{p}^{p} \leq \norm{f}_{p}^{p} + \norm{g}_{p}^{p}$ for $0 < p < 1$. 
    \end{enumerate}
\end{theorem}
\begin{proof}{Proof:}
    For (1), it is clearly enough to take $r = 1$ (replace $f$ and $g$ with $f^{r}$ and $g^{r}$) and $\norm{f}_p = \norm{g}_{q} = 1$ (replace $f$ and $g$ with their normalized versions, and if norm 0, inequality is obvious). The relation $p^{-1} + q^{-1} = 1$ implies $(p - 1)(q - 1) = 1$, so the equations $y = x^{p -1}$ and $x = y^{q - 1}$ are equivalent for $x, y \geq 0$. Thus, $(\cdot)^{p - 1}$ and $(\cdot)^{q - 1}$ are inverses. By Young's inequality, which is depicted below:

    \fig{0.6}{Lecture9Fig1}{}
    we have that,
    \[ \abs{fg} \leq \int_{0}^{\abs{f}}x^{p - 1}dx + \int_{0}^{\abs{g}}y^{q - 1}dy = p^{-1}\abs{f}^{p} + q^{-1}\abs{g}^{q} \]
    Thus, 
    \[ \norm{fg}_{1} \leq p^{-1} \int \abs{f}^{p}d\mu + q^{-1} \int \abs{g}^{q}d\mu = p^{-1} + q^{-1} = 1 \]
    We will prove the second inequality next.
\end{proof}